\chapter{Conclusão}

Este trabalho apresentou o CIC, arquitetura híbrida de orquestração cognitiva para consultas sobre dados de transparência municipal de Caeté (MG), integrando recuperação documental e Text-to-SQL via DuckDB.

A avaliação com conjunto de referência de 80 consultas demonstrou que o subsistema Text-to-SQL com GPT-4.1-mini e esquema minificado alcança elevada taxa de SQL válida (VSR), porém Execution Accuracy modesta sob métrica estrita de equivalência de resultados.

As contribuições incluem: (i)~arquitetura replicável para GovTech municipal; (ii)~\textit{pipeline} ETL e API analítica; (iii)~protocolo de avaliação reprodutível; e (iv)~evidências empíricas para escolha de modelo e tamanho de contexto.

Como trabalhos futuros, propõe-se: revisão humana das referências do conjunto \textit{golden}; avaliação da qualidade das respostas em linguagem natural; testes com usuários reais; e autocorreção iterativa de consultas SQL.
